% Бүлэг 1

\chapter{Онолын үндэс} % Бүлгийн нэр
\label{Chapter1} % Энэ бүлэг рүү ишлэл хийх бол \ref{Chapter1} командыг ашигла 

%-------------------------------------------------------------------------------

% Агуулгад ашигласан хэвшүүлэлтийн зарим командын тодорхойлолт
\newcommand{\keyword}[1]{\textbf{#1}}
\newcommand{\tabhead}[1]{\textbf{#1}}
\newcommand{\code}[1]{\texttt{#1}}
\newcommand{\file}[1]{\texttt{\bfseries#1}}
\newcommand{\option}[1]{\texttt{\itshape#1}}

%-------------------------------------------------------------------------------



%-------------------------------------------------------------------------------

\section{Бүлгийн зорилго ба хамрах хүрээ}

Энэхүү бүлэгт хүйсийн болон нийгмийн тэгш, хүртээмжтэй байдлын тухай үндсэн ойлголт, тэдгээрийн нийгэм дэх ач холбогдол, мөн эдгээр асуудалтай холбоотой мэдээллийг хэрэглэгчдэд ойлгомжтой хэлбэрээр хүргэхэд ашиглагдах хиймэл оюун ухаанд суурилсан чатбот системийн онолын үндсийг авч үзнэ. Үүний хүрээнд хүйсийн тэгш байдал, нийгмийн хүртээмжтэй байдал, ялгаварлан гадуурхалт ба гадуурхлын нийгмийн үр дагавар, байгалийн хэл боловсруулалт(Natural Language Processing), том хэлний загвар(Large Language Model), Retrieval-Augmented Generation (RAG) аргачлал болон хариуцлагатай AI-ийн үндсэн зарчмуудыг тайлбарлана. Мөн эдгээр ойлголтуудыг Монгол Улсын хууль эрх зүйн орчин болон олон улсын зарчимтай уялдуулан авч үзнэ.

%-------------------------------------------------------------------------------

\section{Хүйсийн болон нийгмийн тэгш байдлын онолын үндэс}

\subsection{Хүйсийн тэгш байдлын ойлголт}

Хүйсийн тэгш байдал гэдэг нь хүйсээс үл хамааран хүн бүр эрх, боломж, оролцооны нөхцөлөөр тэгш хангагдахыг хэлнэ. Энэ ойлголт нь зөвхөн хууль эрх зүйн тэгш байдлаар хязгаарлагдахгүй бөгөөд боловсрол, хөдөлмөр эрхлэлт, нийгмийн оролцоо, шийдвэр гаргалт, мэдээлэлд хүрэх боломж зэрэг өргөн хүрээний асуудлыг хамардаг \cite{gender_equality_general}. Хүйсийн тэгш байдлыг хангах нь хувь хүний эрхийг хамгаалахаас гадна нийгмийн шударга байдал, тогтвортой хөгжил, эдийн засгийн оролцоог дэмжих үндэс болдог.

Практикт хүйсийн тэгш бус байдал нь ил тод болон далд хэлбэрээр илэрч болзошгүй. Тухайлбал тодорхой үүрэг, мэргэжил, зан төлөвийг зөвхөн аль нэг хүйст илүү тохирно гэж үзэх хэвшмэл ойлголт нь хувь хүний сонголт, хөгжлийн боломжийг хязгаарлах эрсдэлтэй. Иймээс хүйсийн тэгш байдлыг хангах асуудал нь зөвхөн эрх зүйн зохицуулалт бус, мэдээллийн хүртээмж, боловсрол, нийгмийн ойлголт, харилцааны соёлтой нягт холбоотой.

\subsection{Нийгмийн хүртээмжтэй байдлын ойлголт}

Нийгмийн хүртээмжтэй байдал нь хүн бүр өөрийн гарал үүсэл, хүйс, нас, хөгжлийн бэрхшээл, хэл, соёл, нийгэм эдийн засгийн нөхцөл байдлаас үл хамааран нийгмийн амьдралд тэгш оролцох, үйлчилгээ авах, шийдвэр гаргалтад дуу хоолойгоо хүргэх боломжтой байх нөхцөлийг илэрхийлдэг \cite{social_inclusion_general}. НҮБ-ын Нийгмийн хөгжлийн хөтөлбөрүүдэд нийгмийн оролцоог тэгш хангах, үйлчилгээний тэгш хүртээмжийг дэмжих нь гол зорилго гэж үздэг.

Хүртээмжтэй байдал нь зөвхөн бодит орчны асуудал биш бөгөөд мэдээллийн хүртээмжтэй байдал, ойлгомжтой хэл найруулга, дижитал системд тэгш хандах боломж зэрэгтэй холбоотой. Иймээс нийгмийн хүртээмжтэй орчныг дэмжихэд хэрэглэгч төвтэй, энгийн, ойлгомжтой, ялгаварлан гадуурхахгүй мэдээллийн үйлчилгээ чухал ач холбогдолтой.

\subsection{Ялгаварлан гадуурхалт, дээрэлхэлт ба нийгмийн ач холбогдол}

Хүйсийн болон нийгмийн тэгш бус байдлын нэг илрэл нь ялгаварлан гадуурхалт, гадуурхал, дээрэлхэлт хэлбэрээр илэрдэг. Энэ нь боловсролын орчин, ажлын байр, цахим орчин, нийгмийн өдөр тутмын харилцаанд ажиглагдаж болно. UNESCO-гийн мэдээлснээр дэлхийн хэмжээнд сурагчдын 30 хувиас дээш нь дээрэлхэлтэд өртсөн байдаг бөгөөд энэ нь сурлагын амжилт, сэтгэл зүй, эрүүл мэнд, сургууль завсардалт зэрэгт сөргөөр нөлөөлдөг \cite{school_bullying_unesco}. Иймээс тэгш харилцаа, харилцан хүндлэл, гадуурхлын эсрэг ойлголтыг зөв түгээх нь зөвхөн онолын бус, бодит нийгмийн ач холбогдолтой асуудал юм.

Нөгөө талаас ажлын байр, байгууллагын орчинд тэгш боломж, ялгаварлан гадуурхалтаас ангид байх асуудал нь хөдөлмөрийн эрх, аюулгүй орчин, шударга үнэлгээ, нийгмийн оролцоотой холбогддог. Ийм мэдээллийг зөв эх сурвалжид тулгуурлан, энгийн ойлгомжтой хэлбэрээр түгээх нь хэрэглэгчийн мэдлэгийг нэмэгдүүлэхээс гадна урьдчилан сэргийлэх ач холбогдолтой.

Дэлхийн хэмжээнд хүйсийн тэгш бус байдал, хүчирхийлэл, гадуурхал тулгамдсан хэвээр буйг олон улсын тайлан, статистик мэдээ нотолж байна. Дэлхийн эдийн засгийн чуулганы 2023 оны \textit{Global Gender Gap Report}-д өнөөгийн өсөлтийн хурдаар бүх салбарт хүйсийн тэгш байдалд бүрэн хүрэхэд ойролцоогоор 131 жил шаардагдана гэж тооцоолсон \cite{wef_gender_gap}. Дэлхийн Эрүүл Мэндийн Байгууллагын дүн шинжилгээгээр амьдралынхаа явцад бие махбод эсвэл бэлгийн хүчирхийлэлд өртсөн эмэгтэйчүүд дэлхий нийтийн эмэгтэйчүүдийн 27 орчим хувийг эзэлдэг \cite{who_vaw_2021}, мөн ажиллах хүчний оролцоонд эрэгтэй, эмэгтэйчүүдийн зөрүү ойролцоогоор 25 нэгж хувь хэвээр байгааг НҮБ-ын Эмэгтэйчүүдийн байгууллага мэдээллэсэн байдаг \cite{gender_equality_general, unwomen_facts}. 

Монгол Улсын хувьд хүүхдийн эсрэг хүчирхийллийн нөхцөл байдлын талаарх судалгааг UNICEF Mongolia тогтмол явуулж байгаа \cite{unicef_mongolia_violence}, мөн Үндэсний статистикийн хороо хүйсийн тэгш байдлын үзүүлэлтүүдийг жил бүр нийтэлж байна \cite{nso_mongolia_gender}. Монгол Улсын Хүний эрхийн Үндэсний Комисс жил бүр хүний эрхийн төлөв байдлын талаар албан ёсны илтгэл гаргадаг бөгөөд энэ нь бодлогын болон судалгааны эх сурвалж болдог \cite{nhrcm_report}. Эдгээр үзүүлэлтээс харахад тэгш бус байдал, хүчирхийлэл, гадуурхал нь зөвхөн хувь хүний түвшний асуудал биш, дэлхий нийтийн болон Монгол Улсын хэмжээнд хэвээр буй системчилсэн асуудал юм. Иймээс холбогдох мэдээлэл, эрх зүйн зохицуулалт, тусламжийн сувгуудыг хэрэглэгчдэд хүртээмжтэй хэлбэрээр хүргэх AI шийдэл бий болгох нь нийгмийн хэрэгцээтэй уялдаатай.

%-------------------------------------------------------------------------------

\section{Хиймэл оюун ухаан ба байгалийн хэл боловсруулалт}

\subsection{Хиймэл оюун ухааны үндсэн ойлголт}

Хиймэл оюун ухаан нь өгөгдөлд тулгуурлан суралцах, хэв шинж илрүүлэх, таамаг гаргах, асуудал шийдвэрлэх, хүний хэлтэй төстэй мэдээлэл боловсруулах зэрэг чадварыг компьютерын системд хэрэгжүүлэх технологийн өргөн хүрээний салбар юм. Орчин үеийн AI системүүд нь машин сургалт, гүн сургалт, компьютер хараа, байгалийн хэл боловсруулалт зэрэг олон дэд чиглэлээс бүрддэг.

Нийгмийн чиглэлийн асуудалд AI ашиглах үед зөвхөн техникийн үзүүлэлт бус, өгөгдлийн чанар, тайлбарлах боломж, шударга байдал, хэрэглэгчийн аюулгүй байдал зэрэг хүчин зүйлс давхар чухал болдог. Ялангуяа тэгш байдал, гадуурхал, эмзэг сэдэвтэй холбоотой мэдээллийн системд AI-ийн хариулт зөв, хүндэтгэлтэй, эх сурвалжтай байх шаардлага тавигдана.

\subsection{Байгалийн хэл боловсруулалт}

Байгалийн хэл боловсруулалт (Natural Language Processing, NLP) нь компьютерыг хүний хэлээр бичигдсэн эсвэл илэрхийлэгдсэн мэдээллийг ойлгох, ангилах, боловсруулах, тайлбарлах, хариу үүсгэх боломжтой болгодог салбар юм. NLP-ийн хэрэглээ нь машин орчуулга, асуулт-хариултын систем, текст ангилал, sentiment analysis, chatbot, virtual assistant зэрэг олон төрлийн системд ашиглагддаг.

Чатботын хувьд NLP нь хэрэглэгчийн асуултын утгыг ойлгох, гол түлхүүр санааг ялгах, асуултын зорилгыг тодорхойлох, тохирсон хариултын бүтэц үүсгэхэд чухал үүрэгтэй. Иймээс NLP нь хиймэл оюун ухаанд суурилсан чатботын технологийн гол үндэс болдог.

\subsection{Трансформер архитектур ба том хэлний загвар}

Орчин үеийн NLP-ийн томоохон ахиц нь трансформер архитектурын хөгжлөөс эхэлсэн гэж үздэг. \textit{Attention Is All You Need} өгүүлэлд танилцуулсан трансформер архитектур нь self-attention механизмаар дамжуулан өгүүлбэр дэх үгсийн хоорондын хамаарлыг илүү үр дүнтэй моделчилж, дараалал боловсруулах чанарыг сайжруулсан \cite{vaswani2017}. Үүний дараа BERT зэрэг pretrained language model-ууд гарч ирснээр хэлний контекст ойлгох чанар эрс сайжирсан \cite{bert2019}. Sequence-to-sequence хувирлыг нэгтгэсэн T5 хэв шинж \cite{t5}, цаашлаад GPT-3-аас эхэлсэн few-shot, in-context learning чадвар \cite{gpt3}, мөн хүний санал хүсэлтэд тулгуурлан суралцах (RLHF) аргачлал \cite{instruct_gpt} нь өнөөгийн LLM-ийн үндэс суурийг тогтоосон.

Том хэлний загвар (Large Language Model, LLM) нь их хэмжээний текст өгөгдөл дээр урьдчилан сургагдсан, хэрэглэгчийн асуултанд уян хатан, байгалийн хэлбэртэй хариулт үүсгэх чадвартай загвар юм. Гэвч LLM нь дангаараа ажиллах үед тухайн байгууллага, сэдэв, хууль эрх зүйн тусгай мэдлэгийг бүрэн зөв мэдэхгүй байх, мөн эх сурвалжгүй зохиомол мэдээлэл гаргах эрсдэлтэй. Энэ эрсдэлийг судлаачид stochastic parrots зэрэг шүүмжлэлээр тодорхойлсон бөгөөд \cite{llm_limitations}, галлюцинацийн судалгаанууд галлюцинацийг intrinsic (өгөгдсөн контекстээс хазайх) ба extrinsic (баримтгүй шинэ мэдээлэл нэмэх) гэж ангилдаг \cite{ji_hallucination_survey, hallucination_survey}. Иймээс өндөр нарийвчлал, эх сурвалжийн шаардлагатай орчинд LLM-ийг нэмэлт мэдлэгийн сантай хослуулах шаардлагатай.

Хүснэгт~\ref{tab:llm-landscape}-д орчин үеийн өргөн ашиглагдаж буй LLM-уудын ерөнхий шинж чанарыг танилцуулав. Үнэлгээний шалгуурт олон хэлт чадвар, контекстийн цонхны хэмжээ, нээлттэй эх эсэх, хэрэглэхэд хялбар интерфэйс зэрэг хүчин зүйлс орно.

\begin{table}[H]
\centering
\small
\renewcommand{\arraystretch}{1.2}
\begin{tabular} {|p{2.7cm}|p{2.0cm}|p{1.8cm}|p{1.9cm}|p{4.6cm}|}
\hline
\textbf{Загвар} & \textbf{Компани} & \textbf{Контекст} & \textbf{Нээлттэй} & \textbf{Онцлог} \\ \hline
GPT-3.5 / GPT-4~\cite{gpt3,instruct_gpt} & OpenAI & 8K--128K & Үгүй & RLHF, инструкц дагах өндөр чанартай \\ \hline
Gemini 1.x / 2.x~\cite{gemini_report} & Google DeepMind & 1M хүртэл & Үгүй (API) & Multimodal, маш урт контекст, олон хэлт \\ \hline
Claude 3.x & Anthropic & 200K & Үгүй & Constitutional AI, аюулгүй харилцаа \\ \hline
LLaMA 2 / 3 & Meta & 4K--128K & Тийм (лиценстэй) & Локал орчинд ажиллах боломжтой, open-weight \\ \hline
Mistral / Mixtral & Mistral AI & 32K хүртэл & Тийм (зарим нь) & Mixture-of-Experts, бага зардалтай \\ \hline
T5 / Flan-T5~\cite{t5} & Google & 512--2K & Тийм & Encoder--decoder, fine-tune-д тохиромжтой \\ \hline
\end{tabular}
\caption[Том хэлний загваруудын харьцуулалт]{Том хэлний загваруудын ерөнхий харьцуулалт}
\label{tab:llm-landscape}
\end{table}

\FloatBarrier

Тус судалгаанд Gemini загварыг \cite{gemini_report} сонгосон үндэслэлийг 3-р бүлэгт нарийвчлан тайлбарласан болно: гол шалтгаан нь Монгол хэлний дэмжлэг, олон хэлт чадвар, API-р дамжуулан хурдан холбогдох боломж, контекст цонхны хэмжээ юм.

%-------------------------------------------------------------------------------

\section{Чатбот систем ба баримтад тулгуурласан хариулт үүсгэх аргачлал}

\subsection{Чатбот системийн ойлголт ба ангилал}

Чатбот нь хэрэглэгчтэй текст эсвэл дуу хоолойгоор харилцаж, мэдээлэл өгөх, асуултад хариулах, зөвлөмж хүргэх зориулалттай програм хангамжийн систем юм. Судалгааны бүтээлүүдэд чатботыг дүрэмд суурилсан, retrieval-д суурилсан, генератив, гибрид хэлбэрээр ангилдаг \cite{chatbot_review}. Дүрэмд суурилсан чатбот нь урьдчилан тогтоосон нөхцөл, логик, түлхүүр үг дээр ажилладаг бол генератив чатбот нь хэрэглэгчийн асуултын утгыг илүү өргөн хүрээнд ойлгон шинэ хариулт үүсгэх боломжтой.

Хүйсийн болон нийгмийн тэгш байдалтай холбоотой сэдэвт чатбот нь зөвхөн яриа өрнүүлэхээс гадна хэрэглэгчид мэдээлэл өгөх, зөв ойлголт түгээх, эх сурвалж руу чиглүүлэх, зарим тохиолдолд дэмжих өнгө аястай харилцах шаардлагатай байдаг. Иймд ийм төрлийн систем нь зөвхөн generative байхаас гадна баримтад тулгуурлах чадвартай байх нь чухал.

Чатбот системүүдийг үндсэн ангилалаар нь харьцуулсан үзүүлэлтийг Хүснэгт~\ref{tab:chatbot-types}-д харуулав. Энэхүү харьцуулалт нь зохиомжийн шатанд тохирох төрлийг сонгох үндэс болж өгнө.

\begin{center}
\begin{table}[!htbp]
\begin{tabular}{|p{2.3cm}|p{3.8cm}|p{3.7cm}|p{3.8cm}|}
\hline
\textbf{Төрөл} & \textbf{Гол ажиллах зарчим} & \textbf{Давуу тал} & \textbf{Сул тал} \\ \hline
Дүрэмд суурилсан (rule-based) & Урьдчилан тогтоосон if-then дүрэм, түлхүүр үг, шийдвэрийн мод & Тогтвортой, тайлбарлахад хялбар, эмзэг сэдвийг хатуу хяналтанд авна & Уян хатан биш, NLP-ийн өргөн хүрээний асуултанд хариулж чадахгүй \\ \hline
Retrieval-д суурилсан & Урьдчилан зохиосон хариултын сан, асуултын утга-ойролцоо хайлт & Хариулт нийцтэй, баримтаас хазайхгүй, FAQ-д тохиромжтой & Шинэ агуулга үүсгэж чадахгүй, мэдлэгийн сангаас гадуур асуултанд эмзэг \\ \hline
Генератив (LLM-based) & LLM-ээр хариултыг шинээр үүсгэх~\cite{vaswani2017,gpt3} & Өргөн хүрээний асуултанд уян хатан, хүний хэлбэртэй хариулт & Эх сурвалжгүй, галлюцинацийн эрсдэлтэй~\cite{llm_limitations,hallucination_survey} \\ \hline
Гибрид (RAG) & Retrieval + Generation: мэдлэгийн сангаас баримтыг татаж, LLM-ээр хариу үүсгэх~\cite{rag_survey,grounded_generation} & Эх сурвалжид тулгуурласан, шинэчлэгдэх боломжтой, ишлэлтэй хариулт & Pipeline нийлмэл; retrieval-ийн чанараас хариулт хамаардаг \\ \hline
\end{tabular}
\caption[Чатбот системийн ангилалын харьцуулалт]{Чатбот системийн үндсэн ангилалын харьцуулалт~\cite{chatbot_review}}
\label{tab:chatbot-types}
\end{table}
\end{center}

Энэхүү харьцуулалтаас үзэхэд хүйсийн тэгш байдал, гадуурхлын эсрэг сэдэвт, тогтмол шинэчлэгдэх албан ёсны эх сурвалжид тулгуурлах шаардлагатай орчинд гибрид RAG төрлийн архитектур хамгийн тохиромжтой сонголт болж байна.

\subsection{Чатбот системийн үндсэн бүтэц}

Чатбот систем нь ерөнхийдөө дараах бүрэлдэхүүн хэсгүүдтэй:
\begin{itemize}
    \item хэрэглэгчийн интерфэйс,
    \item асуулт боловсруулах модуль,
    \item мэдлэгийн сан,
    \item retrieval эсвэл search модуль,
    \item хариулт үүсгэх загвар,
    \item хяналт, лог, санал хүсэлтийн модуль.
\end{itemize}

Эдгээр бүрэлдэхүүн хэсгүүд нь харилцан уялдаатай ажилласнаар хэрэглэгчийн асуултыг хүлээн авч, холбогдох мэдээллийг олж, эцсийн хариултыг ойлгомжтой хэлбэрээр хүргэх боломж бүрддэг. Хэрэв систем нь мэдлэгийн сантай холбогдсон бол тухайн сэдэвт илүү бодитой, эх сурвалжтай хариулт өгөх нөхцөл бүрдэнэ.

\subsection{Retrieval-Augmented Generation (RAG)}

Retrieval-Augmented Generation (RAG) нь хэрэглэгчийн асуултад хариу үүсгэхээс өмнө гадны мэдлэгийн сангаас холбогдох мэдээллийг хайж татан авч, түүний дараа том хэлний загвараар хариулт үүсгэдэг аргачлал юм \cite{rag_survey}. Энэ нь генератив чадвар болон search/retrieval чадварыг нэгтгэдэг учраас ялангуяа баримт, бодлого, заавар, хууль эрх зүйн мэдээлэл дээр суурилсан системд тохиромжтой.

RAG-ийн ерөнхий урсгал нь дараах үе шатуудаас бүрдэнэ:
\begin{enumerate}
    \item Баримт бичгийг цуглуулах, цэвэрлэх
    \item Хэсэглэх (chunking)
    \item Embedding үүсгэх
    \item Вектор хайлт хийх
    \item Хамгийн хамааралтай хэсгүүдийг татах
    \item Татсан эх сурвалжид тулгуурлан хариулт үүсгэх
\end{enumerate}

RAG аргачлал нь дараах үндсэн давуу талуудтай:
\begin{itemize}
    \item \textbf{Буруу мэдээлэл үүсэх эрсдэлийг бууруулах:} хариулт нь зөвхөн загварын ерөнхий мэдлэг бус, гадны мэдлэгийн сангаас татсан баримтад тулгуурлана.
    \item \textbf{Мэдлэгийн санг шинэчлэх боломж:} шинэ баримт бичиг, журам, бодлого нэмэгдэхэд загварыг дахин сургахгүйгээр системийн мэдлэгийг өргөтгөж болно.
    \item \textbf{Эх сурвалжийн ил тод байдал:} хариултад ашигласан баримтын хэсгийг ишлэх, тайлбарлах боломжтой.
\end{itemize}

RAG-д вектор хайлт чухал үүрэгтэй. Вектор хайлтын системүүд нь асуулт болон баримтын хэсгийг embedding болгон дүрслээд тэдгээрийн утгын ойролцоо байдлыг тооцоолж хайлт хийдэг. ChromaDB зэрэг embedded вектор сангууд нь ийм хайлтыг metadata filter-тэй хослуулан үр ашигтай хийхэд ашиглагддаг \cite{chromadb}. Retrieval болон generation-ийг ийнхүү нэгтгэх нь open-domain question answering болон knowledge-intensive NLP системүүдэд өндөр ач холбогдолтой \cite{grounded_generation}.

%-------------------------------------------------------------------------------

\section{Хариуцлагатай AI ба хүртээмжтэй харилцааны зарчим}

\subsection{Хариуцлагатай AI-ийн үндсэн ойлголт}

Хүйсийн болон нийгмийн тэгш байдалтай холбоотой сэдэвт AI систем боловсруулахад зөвхөн техникийн үзүүлэлт хангалтгүй бөгөөд ёс зүй, шударга байдал, тайлбарлах боломж, хэрэглэгчийн аюулгүй байдлыг хамтад нь авч үзэх шаардлагатай. Responsible AI-ийн түгээмэл зарчимд fairness, reliability and safety, privacy and security, inclusiveness, transparency, accountability зэрэг ойлголтууд багтдаг \cite{responsible_ai}. Эдгээр зарчим нь NIST-ийн AI Risk Management Framework \cite{nist_ai_rmf} болон Европын Холбооны AI Act-ын \cite{eu_ai_act} ерөнхий чиглэлтэй уялдаатай бөгөөд эмзэг хэрэглээний орчинд AI системийг хариуцлагатай ашиглах үндэс болдог.

Хүснэгт~\ref{tab:rai-principles}-д эдгээр үндсэн зарчмуудыг тэдгээрийн тайлбар болон энэхүү дипломын ажлын системд хэрэгжүүлсэн арга замтай харьцуулан танилцуулав.

\begin{table}[H]
\centering
\footnotesize
\renewcommand{\arraystretch}{1.15}
\begin{tabular}{|p{3.0cm}|p{5.0cm}|p{5.8cm}|}
\hline
\textbf{Зарчим} & \textbf{Тайлбар} & \textbf{Системд хэрэгжүүлэх арга} \\ \hline
Fairness (шударга байдал) & Хэрэглэгчийн хүйс, нас, угсаа, хөгжлийн бэрхшээлээс үл хамаарч тэгш хариу өгөх & Олон эх сурвалжаас баримт татах, төвийг сахисан prompt хэв шинж тогтоох \\ \hline
Reliability and safety (найдвартай, аюулгүй) & Алдаа, эрсдэлийг урьдчилан хянах, аюулгүй харилцаа & Crisis indicator-аар амь насанд аюултай нөхцлийг таних, тусламжийн утас руу шилжүүлэх \\ \hline
Privacy and security (нууцлал) & Хувийн өгөгдлийг хамгаалах, шаардлагагүй хадгалалтаас зайлсхийх & Локал вектор сан, audit log, API руу зөвхөн контекстийн хэсэг дамжуулах \\ \hline
Inclusiveness (хүртээмжтэй байдал) & Олон хэл, олон төрлийн хэрэглэгчид тэгш хүрэх & Монгол хэл дээрх UI, олон хэлт embedding (bge-m3) \\ \hline
Transparency (ил тод байдал) & AI шийдвэрийн үндэслэлийг тайлбарлах, эх сурвалжийг харуулах & Хариултад retrieved sources-г хавсаргах, model\_used талбар лог \\ \hline
Accountability (хариуцлага) & Системийн алдаа, гажуудлын төлөө хариуцагч тал тогтоох & chat\_logs-г хадгалах, санал хүсэлтийн механизм, лог хяналт \\ \hline
\end{tabular}
\caption[Хариуцлагатай AI-ийн зарчмууд]{Хариуцлагатай AI-ийн зарчмууд ба судалгааны системд хэрэгжүүлэх арга~\cite{responsible_ai,nist_ai_rmf}}
\label{tab:rai-principles}
\end{table}
\FloatBarrier

\subsection{Ялгаварлан гадуурхахгүй харилцаа ба хүртээмжтэй дизайн}

Тэгш, хүртээмжтэй байдлыг дэмжих зорилготой чатбот нь хүйс, нас, нийгмийн бүлэг, соёлын ялгаа зэргийг буруу хэвшмэл ойлголтоор илэрхийлэхээс зайлсхийх ёстой. Хэл найруулга нь төвийг сахисан, хүндэтгэлтэй, ойлгомжтой байх нь хэрэглэгчийн итгэлцэлд чухал нөлөөтэй. Мөн хэрэглэгчийг шүүх, буруутгах өнгө аястай бус, мэдээлэл өгөх, чиглүүлэх, ойлголт нэмэгдүүлэх өнгө аястай байх шаардлагатай.

Ийм төрлийн системийн хүртээмжтэй дизайн нь зөвхөн үгийн сонголт бус, хэрэглэгчийн асуултын түвшин, мэдээллийн ойлгомжтой байдал, хариултын бүтэц, эх сурвалж руу шилжих боломж, хэл найруулгын энгийн байдал зэрэгтэй холбоотой. Өөрөөр хэлбэл хүртээмжтэй харилцаа нь интерфэйс, контент, хэл найруулга, алгоритмын хариуцлагыг хамтад нь авч үздэг.

%-------------------------------------------------------------------------------

\section{Хууль эрх зүй ба бодлогын орчин}

\subsection{Монгол Улсын эрх зүйн зохицуулалт}

Хүйсийн болон нийгмийн тэгш байдлыг хангах асуудал нь Монгол Улсын хууль эрх зүйн зохицуулалттай холбоотой. Жендэрийн эрх тэгш байдлыг хангах тухай хууль нь хүйсээр ялгаварлан гадуурхахыг хориглох, бодлого, хөтөлбөр, арга хэмжээнд жендэрийн мэдрэмжтэй хандах үндэс суурийг тогтоодог \cite{mn_gender_law}. Мөн Хөдөлмөрийн тухай хууль нь хөдөлмөр эрхлэлт, ажлын байрны харилцаанд ялгаварлан гадуурхалтаас ангид байх зарчмыг агуулдаг \cite{mn_labor_law}.

Эдгээр хууль эрх зүйн зохицуулалт нь AI чатботын мэдлэгийн санд ашиглах боломжтой баталгаатай эх сурвалж болж өгнө. Ялангуяа хэрэглэгч эрх, үүрэг, ялгаварлан гадуурхалтын талаарх асуулт асуух үед албан ёсны мэдээлэлд тулгуурлан хариулах шаардлагатай.

\subsection{Олон улсын зарчим ба стандарт}

Олон улсын түвшинд ялгаварлан гадуурхалтаас ангид байх, тэгш боломжийг хангахтай холбоотой үндсэн баримт бичгүүд бий. Тухайлбал Олон улсын хөдөлмөрийн байгууллагын 111 дүгээр конвенц нь ажил эрхлэлт, мэргэжлийн орчин дахь ялгаварлан гадуурхалтын эсрэг үндсэн зарчмыг тодорхойлдог \cite{ilo_c111}. Ийм олон улсын баримт бичгүүд нь тэгш байдал, хүртээмжийн тухай ойлголтыг нийгмийн болон байгууллагын түвшинд авч үзэхэд чухал лавлагаа болдог.

%-------------------------------------------------------------------------------

\section{Бүлгийн дүгнэлт}

Энэхүү бүлэгт хүйсийн болон нийгмийн тэгш, хүртээмжтэй байдлын тухай суурь ойлголт, ялгаварлан гадуурхалт ба гадуурхлын нийгмийн ач холбогдол, хиймэл оюун ухаан болон байгалийн хэл боловсруулалтын үндэс, чатбот системийн бүтэц, Retrieval-Augmented Generation аргачлалын үүрэг, мөн хариуцлагатай AI-ийн зарчмуудыг авч үзлээ. Эдгээр онолын үндэс нь хүйсийн болон нийгмийн тэгш, хүртээмжтэй байдлыг хангахад чиглэсэн хиймэл оюун ухаант чатбот системийн архитектур, мэдлэгийн сан, харилцааны зарчмыг тодорхойлоход суурь болох юм.